%% GENERATED by tools/from_docs.py from stepss-docs. Do not edit by hand:
%% edit the page named below in stepss-docs and regenerate.
%% source: stepss-docs/src/content/docs/python/installation.mdx

\chapter{Installing the Python package}\label{doc:py-install}

stepss is a Python interface to the RAMSES dynamic simulator and the Helios AC power-flow engine, enabling users to define test cases, run simulations, extract results, and perform post-processing, all within Python scripts or Jupyter notebooks.

stepss supports Windows, Linux, and macOS, for both dynamic simulation and the bundled \hyperref[doc:py-helios]{Helios power-flow engine}, and integrates with standard scientific Python libraries (NumPy, SciPy, Matplotlib).

\section{Installation}\label{py-install:installation}

Install stepss and all recommended dependencies via pip:

\begin{Verbatim}
pip install jupyter ipython stepss
\end{Verbatim}

Required dependencies (matplotlib, scipy, numpy) are installed automatically.

For a minimal installation (no plotting or notebook support):

\begin{Verbatim}
pip install stepss
\end{Verbatim}

\subsection{Version numbers}\label{py-install:version-numbers}

A stepss version is the bundled RAMSES version, optionally followed by a
counter: version \texttt{3.59} carries RAMSES 3.59. The counter is appended when
stepss is released without a new RAMSES behind it, so a \texttt{3.59.1} would bundle
RAMSES 3.59 as well. A new RAMSES release restarts the sequence with a bare
version.

\texttt{stepss} on PyPI starts at \texttt{3.59}; nothing below that can be installed.

To pin an engine version, pin the matching stepss release:

\begin{Verbatim}
pip install "stepss~=3.59.0"    # any 3.59.x, all bundling RAMSES 3.59
\end{Verbatim}

Check what an installed copy carries:

\begin{Verbatim}
import stepss
print(stepss.__version__, stepss.__ramses_version__, stepss.__helios_version__)
\end{Verbatim}

\section{Linux System Prerequisites}\label{py-install:linux-system-prerequisites}

On Linux, the following system libraries must be installed before running stepss:

\begin{Verbatim}
sudo apt install libopenblas0 libgfortran5 libgomp1
\end{Verbatim}

These packages provide:

\begin{small}\begin{longtable}[]{@{}
  >{\raggedright\arraybackslash}p{(\linewidth - 4\tabcolsep) * \real{0.3258}}
  >{\raggedright\arraybackslash}p{(\linewidth - 4\tabcolsep) * \real{0.6742}}
@{}}
\toprule\noalign{}
\begin{minipage}[b]{\linewidth}\raggedright
Package
\end{minipage} & \begin{minipage}[b]{\linewidth}\raggedright
Purpose
\end{minipage} \\
\midrule\noalign{}
\endhead
\bottomrule\noalign{}
\endlastfoot
\texttt{libopenblas0} & OpenBLAS BLAS/LAPACK routines used by the solver \\
\texttt{libgfortran5} & GNU Fortran runtime required by the Fortran components of RAMSES \\
\texttt{libgomp1} & OpenMP runtime for multi-core parallel execution \\
\end{longtable}\end{small}

On most desktop Linux distributions these are already present. If \texttt{stepss} fails to import with a shared-library error, install the packages above and retry.

\section{macOS System Prerequisites}\label{py-install:macos-system-prerequisites}

On macOS (Apple Silicon), install the GNU compiler runtimes and OpenBLAS via \href{https://brew.sh/}{Homebrew} before running stepss:

\begin{Verbatim}
brew install gcc openblas
\end{Verbatim}

These provide the \texttt{libgfortran}, \texttt{libgomp}, and \texttt{libopenblas} dylibs that the bundled arm64 binaries link against. Intel Macs are not supported.

\section{Run-time observables}\label{py-install:run-time-observables}

\texttt{case.addRunObs(...)} asks the engine to record chosen quantities while a
simulation runs, and RAMSES writes them to a curve file as it goes. Nothing has
to be installed for this.

The file is plain text, with a \texttt{\#} header naming the columns and time in column
one, so any reader takes it:

\begin{Verbatim}
import numpy as np
data = np.loadtxt('temp_display.cur')
\end{Verbatim}

\section{Live plotting}\label{py-install:live-plotting}

\texttt{stepss.monitor} plots chosen quantities \emph{while} the simulation runs, one panel
per observable:

\begin{Verbatim}
import stepss

ram = stepss.sim()
case = stepss.cfg('cmd.txt')
ram.execSim(case, 0.0)                  # initialise, paused at t = 0

mon = stepss.monitor(ram, ['BV 4044', 'MS g6'])
curves = mon.run(step=0.5)              # run to the end of the scenario
\end{Verbatim}

It polls the engine in process and draws with matplotlib, which pip installs
with the package, so nothing further is needed. The
\hyperref[doc:py-api]{API reference} lists the descriptor vocabulary and the
rest of the methods.

\begin{docsnote}{Note}

\texttt{stepss.\_\_runTimeObs\_\_} is always \texttt{True} and is deprecated. Do not branch on it.

\end{docsnote}

\section{Platform Support}\label{py-install:platform-support}

\begin{small}\begin{longtable}[]{@{}
  >{\raggedright\arraybackslash}p{(\linewidth - 6\tabcolsep) * \real{0.2162}}
  >{\raggedright\arraybackslash}p{(\linewidth - 6\tabcolsep) * \real{0.2559}}
  >{\raggedright\arraybackslash}p{(\linewidth - 6\tabcolsep) * \real{0.5279}}
@{}}
\toprule\noalign{}
\begin{minipage}[b]{\linewidth}\raggedright
Platform
\end{minipage} & \begin{minipage}[b]{\linewidth}\raggedright
Library
\end{minipage} & \begin{minipage}[b]{\linewidth}\raggedright
Notes
\end{minipage} \\
\midrule\noalign{}
\endhead
\bottomrule\noalign{}
\endlastfoot
\textbf{Windows} & \texttt{ramses.dll} & Primary platform, full support \\
\textbf{Linux} & \texttt{ramses.so} & Full support (requires the \hyperref[py-install:linux-system-prerequisites]{system libraries above}) \\
\textbf{macOS} & \texttt{ramses.so} & Apple Silicon (arm64) only; requires the \hyperref[py-install:macos-system-prerequisites]{Homebrew runtimes above} \\
\end{longtable}\end{small}

\section{Next Steps}\label{py-install:next-steps}

\begin{itemize}
\tightlist
\item
  \hyperref[doc:py-api]{API Reference}, Full API documentation
\item
  \hyperref[doc:py-examples]{Examples}, Practical usage examples
\end{itemize}
